% ==================================================
% 第五章　結論與未來工作
% ==================================================
\chapter{結論與未來工作}
\label{c:conclusions}

本研究以四輪全向底盤與六軸手臂構成的移動機械臂為對象，
探討動態室內環境中的導航、目標接近與機構操作。
方法結合底盤幾何預測、手臂設定點追蹤模型、
障礙運動估測與視覺目標參考，
在共同控制問題中求解底盤速度與手臂設定點變化率。
本章依第四章的模擬驗證整理研究成果與貢獻，
並就模型移轉、計算期限、接觸處理及感知泛化提出後續方向。

% ==================================================
% 5.1 研究成果與貢獻
% ==================================================
\section{研究成果與貢獻}
\label{sec:research_results_contributions}

本研究由導航控制延伸至底盤與手臂協同操作，
再將視覺估測與任務階段更新接入住宅取物流程。
各項成果的比較條件與主要結果彙整於
\cref{tab:ch5_results_summary}。
導航、獨立抽屜與住宅任務分別採用對應的控制及評估條件，
其數據用於說明各項設計的效果與適用範圍。

\begin{table}[htbp]
    \centering
    \caption{主要研究成果與對應驗證範圍}
    \label{tab:ch5_results_summary}
    \begingroup
    \small
    \renewcommand{\arraystretch}{1.25}
    \begin{tabular}{p{0.16\textwidth} p{0.43\textwidth} p{0.29\textwidth}}
        \hline
        評估項目 & 主要結果 & 驗證範圍 \\
        \hline
        執行模型
        & $50$ ms 預測之最大逐軸 RMSE，
          直接積分為 $25.378$ mrad，增廣模型為 $0.444$ mrad
        & 自由空間驗證資料；增廣模型使用後續實錄設定點 \\
        動態導航
        & 共同 $83$ 條路線全部到達；
          $82$ 條有效淨空路線中有 $1$ 條出現評估外形重疊
        & 指定導航配置；淨空採平面外形近似 \\
        預測時域
        & 五步配置的位置誤差第 $95$ 百分位數為
          $7.27$--$9.75$ mm，單步為 $10.78$--$16.64$ mm
        & 三組抽屜配對；預測長度及成本累加形式共同改變 \\
        操作策略
        & 移動中操作相較停駐後操作，
          完整抽屜任務時間縮短 $13.38$--$14.05\%$
        & 三組配對皆完成既定操作與返回條件 \\
        住宅整合
        & $43$ 次紀錄中有 $29$ 次完成；
          無行人配對完成案例之平均時間相對縮短 $5.42\%$
        & 紀錄包含版本修正；時間比較限八個共同種子的完成案例 \\
        \hline
    \end{tabular}
    \endgroup
\end{table}

% ==================================================
% 5.1.1 幾何預測導航與動態避障
% ==================================================
\subsection{幾何預測導航與動態避障}
\label{subsec:conclusion_geometric_navigation}

導航部分以 SE(2) GMPC 為路徑追蹤基礎，
由障礙追蹤取得幾何與運動估計，
再將 CBF 限制與原始雷射防護分別配置於控制求解及命令執行端。
此設計使追蹤後的障礙資訊能參與預測控制，
並由原始回波補充近距離命令核對。
研究貢獻在於兩個處理層次與底盤執行限制的整合，
以及透過共同路線與消融比較分析其效果。

\cref{sec:dynamic_navigation_comparison}的共同路線比較中，
本研究導航配置、MPPI 與 RPP 分別到達 $83$、$64$ 與 $60$ 條路線。
在 $82$ 條共同有效淨空路線中，
三者的幾何重疊路線數分別為 $1$、$6$ 與 $30$。
另於 $28$ 條消融路線中，
同時啟用 CBF 與雷射防護的配置未記錄到負淨空。
這些結果支持所測配置下的到達與間距表現，
其改善亦伴隨較長的通行時間。
各方法的運動方向與內部控制形式不同，
導航比較中的局部側移候選亦未啟用，
因此上述結果限於當時完整控制鏈的比較。

% ==================================================
% 5.1.2 執行動態建模與全身協同控制
% ==================================================
\subsection{執行動態建模與全身協同控制}
\label{subsec:conclusion_execution_wholebody_control}

針對手臂位置命令介面，
本研究同時保留實際關節位置與位置設定點，
將設定點積分及關節伺服追蹤合成為控制週期的增廣模型。
此模型區分設定點變化率與實際關節速度，
並保留設定點固定後仍會持續產生的追蹤運動。
\cref{sec:execution_model_validation}的離線結果顯示，
在相同驗證起點及已知輸入歷程下，
增廣模型的關節預測誤差低於直接速度積分。
該結果支持命令介面建模的必要性，
線上控制表現則由對應操作實驗另外評估。

在全身控制中，底盤與手臂透過共同末端位姿成本形成耦合，
並將實際關節位置、設定點及輪級命令限制納入最佳化。
\cref{sec:wholebody_horizon_comparison}的三組時域配對中，
五步配置具有較低的整段追蹤誤差、手臂命令增量與反轉率，
同時增加單次控制計算成本。
此結果反映預測時間、終端項位置與成本累加比例的共同作用。
另在三組操作策略配對中，
移動中操作於完成既定操作條件下縮短完整任務時間，
顯示進場、接手與機構運動的共同安排可改善任務效率。

% ==================================================
% 5.1.3 視覺導引與受約束操作整合
% ==================================================
\subsection{視覺導引與受約束操作整合}
\label{subsec:conclusion_visual_constrained_integration}

視覺前端由實例分割與深度量測建立操作幾何，
並使用觀測時刻的機器人狀態完成座標轉換。
操作參考再依可達構型、導航路徑與物件機構生成，
使世界固定目標、門扇圓弧、抽屜滑軌及底盤相對持物目標，
可透過任務階段更新接入共用的全身控制器。
命令執行部分則納入在途命令前推、實際放行回授與停止尾段，
將目標更新及控制段承接放在一致的時間關係下處理。

\cref{sec:residential_integrated_task}呈現開門通行、
圓柱抓取與車尾收納、開抽屜取筆與收納、
關抽屜、回程通行與關門，以及返回起點的整合流程。
在無行人條件的八個共同種子中，
完成案例的移動與停駐策略平均時間分別為 $319.25$ 與 $337.53$ s，
呈現接近與手臂運動重疊所帶來的時間差異。
此住宅停駐策略於接近前停車，後續機構操作仍允許全身協同，
其條件與獨立抽屜的完整停駐對照分別保留。

整合紀錄亦包含未完成與停止核對未通過的案例，
如\cref{sec:performance_failures_limits}所述。
目前證據支持指定模擬條件下的感知、參考與控制銜接，
持物固定相對位姿、幾何近似與模擬配速仍構成適用邊界。
各項貢獻據此界定為執行模型、全身控制與任務整合的設計及驗證，
實體抓持品質與一倍速控制期限則保留為後續工作。

% ==================================================
% 5.2 未來研究方向
% ==================================================
\section{未來研究方向}
\label{sec:future_research_directions}

後續工作將以第四章已辨識的誤差與中止條件為依據，
先補齊預測、放行與實際運動之間的對應，
再逐步移轉至具有感測與接觸不確定性的實體環境。
各階段均需固定控制配置、完成條件及比較區間，
使精度、時間與未完成情形能在共同基準下評估。

% ==================================================
% 5.2.1 實機模型辨識與即時控制驗證
% ==================================================
\subsection{實機模型辨識與即時控制驗證}
\label{subsec:future_physical_identification_realtime}

實機移轉首先需辨識實際位置介面的追蹤動態，
並在不同構型、負載與命令變化下，
檢查一階模型的單步與多步殘差。
若誤差隨運動方向或幅度呈現持續變化，
可再評估速度相關或分段參數模型，
並以時間分離的資料檢查其預測能力。
線上偏置與末端積分補償亦需進行獨立消融，
分別比較預測殘差、末端到位及命令變化，
以釐清模型修正與目標修正的個別作用。

計算執行方面，需完整記錄狀態擷取、求解、發布、放行與生效時間，
補齊未配對及差異提示所對應的命令歷程。
再於固定配置下檢查控制段與停止尾段的逐步預測，
並以實際速度及連續保持條件判定停穩。
一倍速驗證需涵蓋完整感知與控制負載，
除典型耗時外，同時統計尾端延遲、逾時頻率及中止後的運動，
確認 $50$ ms 控制週期與實際承接期限的相容性。

% ==================================================
% 5.2.2 真實抓持與接觸動態處理
% ==================================================
\subsection{真實抓持與接觸動態處理}
\label{subsec:future_grasp_contact_dynamics}

實體物件操作需進一步描述摩擦、接觸位置與相對滑移，
使固定相對位姿假設能由可量測的抓持條件逐步取代。
後續可結合夾爪位置、接觸力及物件相對運動，
辨識抓持建立、穩定搬運與失去抓持等狀態，
再評估力或柔順控制對門扇及抽屜操作的作用。
筆狀代理物亦需擴展至不同直徑與材質，
分析物件尺寸及接觸條件對可用接近方向與夾持能力的影響。

幾何預測需將持物、手臂掃掠範圍與盒內既有物件納入檢查，
並檢驗感測誤差與離散取樣對最小淨空的影響。
機構操作的停止過程亦須保留接觸耦合與關節追蹤，
針對尾段結束後仍存在的殘餘運動，
調整停止時長、命令安排與機構相容條件，
再由獨立量測核對實際停止結果。

% ==================================================
% 5.2.3 視覺感知與任務泛化能力
% ==================================================
\subsection{視覺感知與任務泛化能力}
\label{subsec:future_perception_task_generalization}

感知部分需加入實體 RGB-D 資料、相機標定誤差、
深度缺值、遮蔽及定位漂移，
分別評估分割品質、三維目標誤差與操作完成情形。
目標估測除提供位置與方向外，
亦可將資料年齡及幾何不確定性傳遞至選位與接近規則，
使重新觀測、降低速度或停止推進能依量測品質調整。
對家具尺寸及機構軸向的先驗，
則逐步改由多視角觀測與線上估測建立，
以擴大方法可處理的物件及場景範圍。

動態場景驗證需固定障礙運動規則，
分別比較原導航參考、速度調整及速度與側移共同調整的配置，
並檢查避讓後的手臂可達性、目標可見性與任務進度。
後續亦需增加共同條件下的重複測試，
同時呈現完成、未完成及停止結果，
並將完成案例的條件時間與整體任務可靠性分別評估。
透過這些驗證，進一步釐清全身協同運動的時間收益，
以及其在感知、接觸與動態通行條件改變時的適用範圍。
